\documentclass[11pt]{article}
\usepackage[utf8]{inputenc}\usepackage[T1]{fontenc}
\usepackage[a4paper,margin=2.2cm]{geometry}
\usepackage{amsmath,amssymb}
\usepackage{booktabs,array}
\usepackage{xcolor}
\usepackage{enumitem}
\usepackage[colorlinks=true,linkcolor=black,urlcolor=blue]{hyperref}
\usepackage{titlesec}
\titleformat{\section}{\large\bfseries\color{black!80}}{\thesection.}{0.6em}{}
\setlength{\parindent}{0pt}\setlength{\parskip}{5pt}
\newcommand{\code}[1]{\texttt{\small #1}}
\newenvironment{quotebox}{\begin{center}\begin{minipage}{0.92\textwidth}\itshape\color{black!75}}{\end{minipage}\end{center}}

\begin{document}

\begin{center}
  {\LARGE\bfseries The HyMeKo System-Engineering View}\\[3pt]
  {\large A manifesto --- one source, many views}\\[6pt]
  {\small Companion to \emph{HyMeKo Capability Evidence Ledger}. Compiled 2026-06-20.}
\end{center}
\vspace{4pt}

\section{Thesis}
\begin{quotebox}
\textbf{One source, many views.} A robotic system is described \textbf{once}, as a signed
hypergraph in HyMeKo. Every artifact a downstream tool needs --- URDF, SDF, MJCF, a Gazebo launch
bundle, a SysML block diagram, a Lean4 proof obligation, a DOT graph, the figure in a paper, an RL
\code{AgentSpec} --- is a \textbf{projection} (a \emph{view}) of that single source, never a parallel
hand-maintained copy.
\end{quotebox}

This is \textbf{Model-Based Systems Engineering (MBSE)} taken literally. In INCOSE's terms, MBSE is
\emph{the formalized application of modelling to support system requirements, design, analysis,
verification and validation, beginning in the conceptual design phase and continuing throughout
development and later life-cycle phases} (INCOSE SE Vision; the direction of travel of the discipline
through SE Vision 2035). HyMeKo is the \textbf{authoritative source of truth} that vision calls for:
the model is not documentation \emph{about} the system; the model \textbf{is} the system of record.
Code, simulation assets, and diagrams are generated downstream and are, by construction, consistent
with each other because they share one upstream cause --- a single point on the \textbf{digital
thread}.

The failure mode this exists to kill is the \emph{document-centric} trap: the same robot encoded
three times --- once in a URDF, once in a slide, once in a training script --- drifting out of sync
the moment anyone edits one of them. Under MBSE there is nothing to keep in sync, because there is
only one editable artifact and everything else is derived from it.

\section{Two layers, separated on purpose}
Every HyMeKo description separates \textbf{vocabulary} from \textbf{instance}.

\begin{center}\small
\begin{tabular}{@{}p{2.6cm}p{7.2cm}p{4.4cm}@{}}
\toprule
\textbf{Layer} & \textbf{What it is} & \textbf{Example} \\
\midrule
Vocabulary (\code{meta\_*}) & The \emph{kinds} --- link, joint, axis, sensor, context, signal,
  aggregation. Declared once per ecosystem. & \code{meta\_kinematics.hymeko}, \code{meta\_context.hymeko} \\
\addlinespace
Instance & The \emph{specific} robot --- its links, joints, context graph. Imports a vocabulary;
  declares only what is unique. & \code{wam.hymeko}, \code{hymeko\_robot\_reuse.hymeko} \\
\bottomrule
\end{tabular}
\end{center}

A vocabulary is amortised across \textbf{every} system that imports it. The cost of defining ``what a
revolute joint is'' is paid once, not once per robot.

\section{Why this is the selling point, not a refactor}
Measured (\code{reports/2026-06-20-mdsd-reuse-scenario.md}). The ROS2 demo robot, authored both
ways --- bare inline types (kept as a frozen control) vs.\ imported vocabulary:
\begin{itemize}[nosep]
  \item \textbf{Length:} $189\to126$ lines ($-33\%$); $108\to76$ code lines for \emph{one} robot.
        The shared vocabulary is paid once, so the marginal saving for the $N$-th robot is the full
        re-declaration.
  \item \textbf{Semantics, not just brevity:} the imported form carries \emph{typed} joints with
        real axis vectors and limits; URDF/SDF emit valid \code{<joint>} elements. The bare form's
        generic \code{joint} emits \textbf{zero}. Reuse is the difference between a description that
        \emph{projects to a working robot} and one that does not.
\end{itemize}
The length number is the headline; the semantic number is the argument.

\section{What counts as a ``view'' (ISO/IEC/IEEE 42010)}
\emph{View} is used in the precise sense of \textbf{ISO/IEC/IEEE 42010}, the architecture-description
standard INCOSE builds on: a \textbf{view} addresses a set of \textbf{stakeholder concerns} from a
\textbf{viewpoint}. The architecture description \emph{is} the signed hypergraph; each emitter is a
viewpoint; each emitted artifact is a view.

\begin{center}\small
\begin{tabular}{@{}p{5.6cm}p{4.0cm}p{4.6cm}@{}}
\toprule
\textbf{Stakeholder concern} & \textbf{Viewpoint (emitter)} & \textbf{View (artifact)} \\
\midrule
``Will it move correctly in sim?'' & URDF / SDF / MJCF / Gazebo & simulation scene \\
``How does it decompose into blocks \& ports?'' & SysML & block / IBD diagram \\
``Is this property provable?'' & Lean4 & proof obligation \\
``What is the structure?'' & DOT / Mermaid / WASM editor & graph \\
``What is the learning problem?'' & RL \code{AgentSpec} & obs + action + reward + vertices \\
``Did the projection lose information?'' & HyMeKo $\to$ HyMeKo & canonical round-trip \\
\bottomrule
\end{tabular}
\end{center}

The views are \textbf{not} independent documents to be reconciled --- they are consistent by
construction. Adding a view means writing one viewpoint against the IR, never asking the author to
maintain another copy.

\section{The boundary rules (so the principle does not rot)}
\begin{enumerate}[nosep]
  \item \textbf{Vocabulary is imported, never re-declared.} Inline kinds where a \code{meta\_*}
        profile already declares them is duplication in the costume of self-containment. The only
        exception is a \textbf{frozen reference artifact}, marked as such.
  \item \textbf{Parametric difference $\to$ configuration; structural difference $\to$ a different
        description.} Different link lengths = same description, new field values. Different topology
        = a different instance graph. Never a runtime \code{if}.
  \item \textbf{A view is read-only downstream.} Never hand-edit a generated URDF. If a view is
        wrong, the bug is in the source or the emitter.
  \item \textbf{One source per system, even across tools.} The robot that drives the MuJoCo env, the
        Gazebo demo, and the SysML diagram is \emph{one} \code{.hymeko}.
\end{enumerate}

\section{Alignment with INCOSE systems-engineering practice}
This manifesto is the project's reading of established SE doctrine, applied to a hypergraph core.

\begin{center}\small
\begin{tabular}{@{}p{5.4cm}p{6.0cm}p{3.6cm}@{}}
\toprule
\textbf{HyMeKo} & \textbf{Systems-engineering concept} & \textbf{Source} \\
\midrule
The \code{.hymeko} source & Authoritative source of truth / single system model & INCOSE MBSE, SE Vision 2035 \\
Emitter + emitted artifact & Viewpoint + view addressing a concern & ISO/IEC/IEEE 42010 \\
Vocabulary vs instance & Separation of concerns; reference vs system data & ISO/IEC/IEEE 15288 \\
Cross-file import + canonical hash & Traceability; configuration identity & INCOSE Handbook \\
\code{.hymeko} $\to$ URDF/SysML/Lean4/\dots & The digital thread across life-cycle activities & Digital Engineering \\
Round-trip + validation-gated emit & Verification \& Validation & ISO/IEC/IEEE 15288 \\
\bottomrule
\end{tabular}
\end{center}

\textbf{Where this sits on the Vee.} The SE \emph{Vee} has a left leg (decomposition \& definition:
needs $\to$ requirements $\to$ architecture $\to$ design) and a right leg (integration \&
verification). HyMeKo's authoring sits on the \textbf{left leg}, captured once as the hypergraph; the
emitters carry that single definition rightward into each V\&V activity --- the same model that
\emph{defines} the robot is the one \emph{verified} against each consumer.

\textbf{Verification vs validation, kept distinct (INCOSE's two questions).}
\begin{itemize}[nosep]
  \item \emph{Verification --- ``did we build the system right?''} Are the views \textbf{faithful} to
        the source? Mechanically checkable: the HyMeKo$\to$HyMeKo round-trip and the canonical hash
        are verification gates. The MDSD reuse result is a verification statement --- the imported
        description projects to the same kinematic structure, hash-confirmed.
  \item \emph{Validation --- ``did we build the right system?''} Does the model meet the need? This
        the framework \textbf{cannot} assert for you; it requires running the robot, the sim, the
        experiment. The capability ledger is deliberately honest about this line: \textsc{proven}
        rows are largely verification; \textsc{shipped} / \textsc{research} rows are where validation
        is still open.
\end{itemize}

\textbf{Why a hypergraph, in SE terms.} SysML models systems as blocks with ports and pairwise
connections. Many real engineering relations are \textbf{not} pairwise --- a constraint over three
quantities, a context fusing five signals, a signed balance over a cycle. Forcing those into binary
connectors loses information at authoring time. The signed hypergraph (GGK 4-tuple
$K=(B,G,\mu,r)$) keeps higher-order, signed, cross-context relations as first-class, and the pairwise
SysML view is recovered by projection --- never the other way around.

\section{The standing invariant}
\begin{quotebox}
If a fact about a system is true, it is stated \textbf{once} in HyMeKo and \textbf{derived everywhere
else}. Any place a fact is restated by hand is a latent inconsistency, and is treated as a defect ---
not a convenience.
\end{quotebox}
This is the single-authoritative-source principle of MBSE stated as an operational rule, and it is
what makes traceability mechanical rather than aspirational: every view traces to exactly one source
element, and the canonical hash detects when that link is broken.

\end{document}
